\chapter{Security and trust}
\label{ch:security}

\begin{table}[H]
\centering
\begin{tabularx}{\textwidth}{L{0.26\textwidth} Y}
\toprule
\textbf{Item} & \textbf{Decision} \\
\midrule
Authentication & IDTOKENS; one token per user in \file{\textasciitilde/.condor/tokens.d/} \\
Transport & A single TCP port, 9618, via \kn{condor\_shared\_port} \\
Entry-node accounts & System accounts only --- \kn{--no-create-home --shell /usr/sbin/nologin} \\
Account per user & Distinct UIDs, never a shared one --- required for fair-share and spool isolation \\
Worker accounts & None, ever. Jobs run as \kn{nobody} in a scratch sandbox \\
Job isolation & \textbf{Undecided} --- see \S\ref{sec:isolation} \\
\bottomrule
\end{tabularx}
\caption{Security decisions.}
\end{table}

\section{Why tokens rather than a filesystem trust model}

The design goal that users never log into the entry node removes the usual basis
for trust. There is no shared home directory to hold an authorized\_keys file and
no shared filesystem whose UIDs mean anything.

IDTOKENS fits exactly: a token is a bearer credential tied to a principal in the
\kn{UID\_DOMAIN} namespace, independent of any Unix account on any machine.

\begin{keypoint}
The single port matters as much as the authentication. \kn{condor\_shared\_port}
funnels every daemon conversation through TCP 9618, which is what makes the whole
pool reachable through one ssh tunnel --- and therefore what makes ``submit from
your own PC'' practical rather than theoretical.
\end{keypoint}

\section{User tokens and daemon tokens are different}

This distinction costs an hour if missed, because the failure is silent at the
point where you are looking.

\begin{trap}
A \textbf{user} token does not satisfy \kn{ALLOW\_DAEMON}. Without a daemon
credential a worker starts cleanly, fails to register, and simply never appears
in \kn{condor\_status}. No error surfaces at the client.
\end{trap}

Daemon tokens live in \file{/etc/condor/tokens.d/}, owned by root, mode 600. User
tokens live in the user's own \file{\textasciitilde/.condor/tokens.d/}.

\begin{trap}
Do not test daemon authentication with \kn{condor\_ping ... DAEMON} as an
unprivileged user. It cannot read the root-owned token, falls back to SSL, and
blocks on an interactive trust prompt --- which looks like a hang, not like a
permissions problem. Check pool membership instead.
\end{trap}

\section{Job isolation is the open decision}
\label{sec:isolation}

As built, jobs run as \kn{nobody} in a scratch sandbox. That bounds what they can
write, and the sandbox is removed as soon as output transfer completes.

It does not bound what they can \emph{read}. A job can read world-readable files
and use the network on any machine it lands on --- and in this pool, the machines
it lands on are colleagues' personal desktops.

\begin{keypoint}
\textbf{This must be decided before the pool grows, not after.} If bare execution
is unacceptable, the answer is Apptainer for every job from day one.
Retrofitting containerisation onto an established pool is far harder than
starting with it, because every existing submit file becomes wrong at once.
\end{keypoint}

The tension is real and worth stating plainly: containers cost submission
ergonomics, and submission ergonomics is the entire reason this pool exists.
A design that makes submitting harder than running locally will not be used, and
an unused pool has no security problems at all.
